\clearpage
\subsection{自同构}\label{sec:12i}
\sr{sec:10}{第~10~节}关于 $\F(M)^{\mathrm{pr}}$ 的全部观察都可推广到 $\F(M)$。当然，有些陈述须略作改变，以容纳膜结构，但通常很直接。作为实际演示，考察一个较有吸引力的例子：辛自同构群在 Fukaya 范畴上的作用，最初在\secref{10b}讨论。

先说一些一般事实。取任意辛流形 $(M,\allowbreak \omega)$，配有殆复结构 $I_M$ 和二次复体积形式 $\eta_M^2$。
\begin{definition*}
$M$ 的\emph{分次自同构}是二元组 $(\phi,\allowbreak \gamma^\#)$，由辛自同构 $\phi$ 和映射 $\gamma^\#:\operatorname{Gr}(TM)\to\R$ 组成，对所有 $\Lambda\in\operatorname{Gr}(TM)$ 满足
\[
\exp(2\pi i\gamma^\#(\Lambda))=\frac{\alpha_M(D\phi(\Lambda))}{\alpha_M(\Lambda)}.
\]
\end{definition*}
\SourcePage{196}
这种自同构的复合定义为 $\phi_2^\#\circ\phi_1^\#=(\phi_2\circ\phi_1,\allowbreak \gamma_2^\#\circ D\phi_1+\gamma_1^\#)$。分次自同构自然作用于 Lagrangian 膜的集合：若 $(L,\allowbreak \alpha^\#,\allowbreak P^\#)$ 为膜，则像 $\phi(L)$ 继承分次 $x\mapsto\gamma^\#(TL_{\phi^{-1}(x)})+\alpha^\#(\phi^{-1}(x))$，以及推前 Pin 结构 $(\phi^{-1})^*P^\#$。这还不是最一般的情形。
\begin{definition*}
\emph{膜自同构}是三元组 $\phi^\#=(\phi,\allowbreak \gamma^\#,\allowbreak \beta)$，其中 $(\phi,\allowbreak \gamma^\#)$ 如前，$\beta\to M$ 为实线丛。
\end{definition*}
这种自同构在膜上的作用还包括 Pin 结构的扭曲：
\begin{equation}\label{eq:12.29}
P^\#\longmapsto(\phi^{-1})^*(P^\#\otimes\beta|L).
\end{equation}
一个相关例子是平移运算 $S$；由\stref{笔记}{11.18}，它可写为 $(\Id_M,\allowbreak \gamma^\#=-1,\allowbreak \beta=\xi_M)$ 的作用。

以下假设 $M$ 为带角正合辛流形，考虑由底层 $\phi$ 属于 $\operatorname{Aut}(M,\allowbreak \partial M)$ 的 $(\phi,\allowbreak \gamma^\#)$ 组成的群 $\operatorname{Aut}^{\mathrm{gr}}(M,\allowbreak \partial M)$。赋予它从 $\operatorname{Aut}(M,\allowbreak \partial M)\times C^\infty(\operatorname{Gr}(TM),\allowbreak \R)$ 诱导的拓扑。$\operatorname{Aut}^{\mathrm{gr}}(M,\allowbreak \partial M)$ 在 $H\F(M)$ 上有典范严格作用，并下降为 $\pi_0(\operatorname{Aut}^{\mathrm{gr}}(M,\allowbreak \partial M))$ 的弱作用。对 Fukaya 范畴本身，使用小节 \sr{sec:10b}{(10b)}、\sr{sec:10c}{(10c)} 的同一构造，得到 $\operatorname{Aut}^{\mathrm{gr}}(M,\allowbreak \partial M)$ 在 $\F(M)$ 上的协调作用，同样下降为 $\pi_0(\operatorname{Aut}^{\mathrm{gr}}(M,\allowbreak \partial M))$ 的弱作用。现在加入一般膜自同构，但不区分 $M$ 上同构的线丛。这定义 $\operatorname{Aut}^{\mathrm{gr}}(M,\allowbreak \partial M)$ 关于 $H^1(M;\Z/2)$ 的扩张，记为 $\operatorname{Aut}^\#(M,\allowbreak \partial M)$。该群在 Fukaya 范畴、甚至 Donaldson–Fukaya 范畴上的自然作用，一般都只有弱作用，因为\eref{12.29}要求在同构类 $[\beta]$ 中选取代表 $\beta$。这一现象在实际应用中并不十分重要，不过后面讨论对合时，细心的读者仍会看到它的痕迹，见\sr{sec:14}{第~14~节}。\znote{188}{“弱作用”指群作用只在同构意义下成立：公理 $g(hx)\cong(gh)x$、$ex\cong x$ 只要求自然同构而非等式。膜自同构多出的线丛 $\beta$ 使 (12.29) 必须先选代表，故作用只能弱；把线丛自身的对称群 $H^0(M;\Z/2)$ 一并记入，$\operatorname{Aut}^\#(M,\partial M)$ 就成为弱 $2$ 群——群公理被范畴化，元素之间的同构本身也成为数据。也正因如此，只有 $\pi_0$ 才是通常意义上的群。笔记 12.5 的计算说明，平凡线丛的自同构 $-\Id_\beta$ 给出的正是自然变换 $-E$。[来源：Seidel, Fukaya categories and Picard–Lefschetz theory, §10b–10c 与笔记 12.5。]}
\begin{note}\label{st:12.5}
为说明代表 $\beta$ 的选择确有影响，考虑三个正合 Lagrangian 膜 $L_k^\#=(L_k,\allowbreak \alpha_k^\#,\allowbreak P_k^\#)$（$k=0,\allowbreak 1,\allowbreak 2$），其中后两个同构：具有相同的 Lagrangian 子流形及分次，Pin 结构同构但不一定严格相等。Pin 同构 $g:P_1^\#\to P_2^\#$ 的选择确定 Floer 上同调群同构
\begin{equation}\label{eq:12.30}
HF^*(L_0^\#,L_1^\#)\cong HF^*(L_0^\#,L_2^\#).
\end{equation}
为简单起见，假定 $L_0$ 与 $L_1=L_2$ 横截相交，这样相关 Floer 数据中的 Hamilton 项可取为零。固定交点 $y$，令 $\Lambda_k^\#=\Lambda_{k,\allowbreak y}^\#$（$k=0,\allowbreak 1,\allowbreak 2$）为由 $L_k^\#$ 在该点得到的线性膜。

选择处于正确同伦类中的路径 $\lambda\in\Omega^-(\operatorname{Gr}(TM_y),\allowbreak \Lambda_0,\allowbreak \Lambda_1)$，以及相应向量丛 $F_\lambda\to[0;1]$ 上的 Pin 结构 $P_\lambda^\#$。可将
\SourcePage{197}
$(\delta_H)_\lambda$ 同时与定向空间 $o(\Lambda_0^\#,\allowbreak \Lambda_1^\#)$、$o(\Lambda_0^\#,\allowbreak \Lambda_2^\#)$ 识别，见\stref{引理}{11.11}。但这种识别依赖辅助数据的选择，即同构 $f_0:(P_\lambda^\#)_0\to P_0^\#$ 和 $f_k:(P_\lambda^\#)_1\to P_k^\#$、$k=1,\allowbreak 2$。为得到\eref{12.30}底层的链映射，应取 $f_2=g_y\circ f_1$，并考虑诱导的定向空间映射。反转 $g$ 的符号会改变 $f_2$ 的符号，因而按\stref{笔记}{11.19}的讨论，改变识别 $(\delta_H)_\lambda\cong o(\Lambda_0^\#,\allowbreak \Lambda_2^\#)$，最终将\eref{12.30}乘以 $-1$。
\end{note}

改善这一情形的一种办法，是把 $\operatorname{Aut}^\#(M,\allowbreak \partial M)$ 看作更精细的对象——“弱 $2$ 群”，将线丛的内部对称群 $H^0(M;\Z/2)$ 也纳入考虑。这与一般范畴论十分吻合，因为范畴论本就认为“一个范畴的自同构群”是这类对象。用这一语言，\stref{笔记}{12.5}中的计算说明：平凡线丛 $\beta=M\times\R$ 的自同构 $-\Id_\beta$ 给出恒等函子到自身的自然变换 $-E$。

\subsection{补充说明}\label{sec:12j}
Fukaya 范畴还有许多其他性质，这里无法讨论。不过，其中两项值得作最基本的介绍，至少因为它们在上同调层次的对应物已在 Floer 理论的讨论中出现。

第一项是 PSS 同构 $HF^*(L^\#,\allowbreak L^\#)\cong H^*(L;\K)$。相应的链层次陈述说：若 $L^\#$ 是单个正合 Lagrangian 膜，则只考察这个对象得到的 $\Ai$ 代数 $\B=\operatorname{hom}_{\F(M)}(L^\#,\allowbreak L^\#)$，拟同构于用 $L$ 的经典拓扑所得到的代数，例如奇异上链的微分分次代数，或等价地，使用 Morse 同调及 Fukaya 的梯度树运算\bcite{54}得到的代数。该陈述的早期形式见\bcite{58}，其中通过构造模空间的同构，基本上旨在得到相关代数结构的相等。另一种办法是在\bcite{33}及其前期工作的方式下，混合使用伪全纯多边形与梯度树；它的分析技术较不棘手，但只能证明拟同构。还有允许扭曲系数的小变体，其中包括\eref{12.14}的链层次版本。

第二项是 Donaldson–Fukaya 范畴的 Frobenius 性质的对应物，所以假设 $\xi_M$ 平凡。对任意 $\Ai$ 范畴 $\A$，有典范 $\Ai$ 函子 $\mathcal D:\A\to\Q=\modcat(\A)$，称为\emph{抽象 Serre 函子}。对象 $Y$ 在 $\mathcal D$ 下的像为对偶 $\Ai$ 模 $\Y^\vee$，满足 $\Y^\vee(X)=\operatorname{hom}_{\A}(Y,\allowbreak X)^\vee$，结构映射也有类似公式。\znote{189}{Serre 函子的定义性性质是自然同构 $\operatorname{hom}_{\A}(A,B)\cong\operatorname{hom}_{\A}(B,SA)^\vee$；在凝聚层的导出范畴中它由 $(\cdot\otimes\omega_X)[n]$ 实现，因此“$n$ 维弱 Calabi–Yau”就是把 Serre 函子识别为平移 $[n]$，即 $\mathcal D\cong\mathcal I[n]$。这里 $\mathcal D$ 把对象 $Y$ 送到对偶模 $Y^\vee$，正是把 12e 中带 $\lambda^{\mathrm{top}}(TL)$ 扭转的 Poincaré 对偶配对搬到 $\Ai$ 层面，所以弱 CY 性质可读作“Frobenius 性质差一个 $[n]$”。把 $\A$ 换成 $\F(M)$ 即得 Donaldson–Fukaya 范畴，而 $\mathcal D$ 与 Yoneda 嵌入 $\mathcal I$ 的差别正是那个扭转在模层面的体现，也是 12j 全节的核心。[来源：nLab “Serre functor”，https:/\allowbreak /\allowbreak ncatlab.org/\allowbreak nlab/\allowbreak show/\allowbreak Serre+\allowbreak functor（2026-11-15 访问）；并见原书 §12j。]}
\begin{definition*}
若在 $H(\operatorname{fun}(\A,\allowbreak \Q))$ 中，存在从 Yoneda 嵌入 $\mathcal I$ 到 $\mathcal D$ 的次数为 $-n$ 的
\SourcePage{198}
同构，则称 $\A$ 为 $n$ 维\emph{弱 Calabi–Yau 范畴}。
\end{definition*}
自然变换 $\mathcal I\to\mathcal D$ 自然来自对偶 Hochschild 上同调 $HH(\A,\allowbreak \A^\vee)$ 中的类。在我们的情形 $\A=\F(M)$ 中，相关类通过考察只有出向无穷远点、并额外带一个内部标记点的带标记点圆盘来定义，方式如\secref{10c}。随后由\stref{引理}{1.6}及 Floer 上同调中的 Poincaré 对偶，得到相应自然变换为同构。
